% TOP_PROBLEM: {"problemId": "problem.positivity-problem-for-linear-recurrence-sequences", "canonicalTitle": "Positivity Problem for Linear Recurrence Sequences", "releaseRank": 296, "primaryDomain": "logic_foundations_set_theory", "primaryDomainLabel": "Logic, foundations & set theory", "displayStatus": "open_with_solved_subcases", "releaseStatus": "open_with_solved_subcases"}
% !TeX program = lualatex
% Definition and sources only; this file is not an LLM solution attempt.
\documentclass[11pt]{article}
\usepackage[margin=25mm]{geometry}
\usepackage{fontspec}
\setmainfont{Latin Modern Roman}
\usepackage{amsmath,amssymb,mathtools}
\usepackage{unicode-math}
\setmathfont{Latin Modern Math}
\usepackage{xurl,hyperref}
\hypersetup{colorlinks=true,urlcolor=blue}
\setcounter{secnumdepth}{0}
\setlength{\parindent}{0pt}
\setlength{\parskip}{5pt}
\setlength{\emergencystretch}{3em}
\begin{document}
\section*{\texorpdfstring{Positivity Problem for Linear Recurrence Sequences}{Positivity Problem for Linear Recurrence Sequences}}\label{296}
\subsection{Definitions and mathematical statement}
Input consists of $k\ge1$, rational numbers $a_1,\ldots,a_k$, and rational initial values, all exactly encoded in binary, defining
\[
 u_{n+k}=\sum_{i=1}^k a_i u_{n+k-i}\qquad(n\ge0).
\]
The target is a single total algorithm deciding whether $u_n\ge0$ for \emph{every} $n\ge0$. No running-time bound is imposed, and $k$ is not fixed. Coefficients may vanish, including $a_k$, so the displayed recurrence need not have minimal order $k$. Eventual nonnegativity or nonnegativity up to a supplied finite index is a different decision problem. Equality to zero is permitted in the desired sequence.
\par\smallskip\textit{Formulation and concept references:} \href{https://arxiv.org/pdf/2302.13675v2}{[S1]}.

\subsection{Short English statement}
Can an algorithm always determine whether every term of an arbitrary rational linear recurrence sequence is nonnegative?
\subsection{Sources}
\begin{itemize}
\item[S1]Positivity-hardness results on Markov decision processes, arXiv2302.13675v2, 5 April 2024; TheoretiCS3 Article9.
\url{https://arxiv.org/pdf/2302.13675v2}.
\textit{Formulation source.}
\end{itemize}
\end{document}
